\subsection{Setup and relation to the supplied sources}

Let $F$ be any field and let $B=\{u_b:b\in I\}$ be a nonempty finite set
of distinct elements of $XF[[X]]$, with $|I|=N$. Put
\[
 v=\ord_X,\qquad s_{bc}=v(u_b-u_c)\quad(b\ne c),\qquad s_{bb}=+\infty.
\]
The finite contacts are positive integers and obey the ultrametric inequality.
For a precision vector $\mathbf n=(n_b)_{b\in I}\in\mathbb Z_{\ge0}^I$ define
\[
 \Omega_{\mathbf n}(H)_b=H(X,u_b(X))\bmod X^{n_b},\qquad
 J_d(H)=H\bmod (X,Y)^d.
\]
A record with $n_b=0$ is empty. Throughout the target order $d$ is positive.
The higher coefficients of $H\in F[[X,Y]]$ are arbitrary. Thus this is
recovery of a quotient of the source, rather than ordinary decoding of a
bounded-degree polynomial.

The supplied article~\cite{Source} proves the common-length observation kernel and distance
formula in the results labelled \texttt{lub:thm:kernel} and
\texttt{lub:thm:distance}; it gives a common-length tree recursion in
\texttt{lub:thm:dp}, and the least common precision for maximum distance in
\texttt{lub:thm:saturation}. Its question \texttt{lub:q:adversarial} asks for
unequal lengths, higher coefficient precision, and faster error-location
algorithms. For exactly $d$ branches, unequal-length clean recovery is already
covered by \texttt{lub:prop:threshold}. The present extension concerns any
number of redundant branches, arbitrary unequal lengths, and errors or
erasures. Its local setting is also obtained by reducing the Lubin--Tate
branches used by the Lean development modulo the uniformizer; there
\[
 s_{bc}=q^{v_\pi(b-c)}.
\]
No Riemann hypothesis is an assumption of any result in this section.

The immutable source versions and inspection scope are recorded in Section~\ref{sec:sources}.

\subsection{Replacing every ambiguity by branch factors}

The central observation eliminates the unknown coefficients from the
distance problem. A coefficient unit is immaterial to valuations.

\begin{lemma}[Product domination for polynomials]\label{lem:domination}
Let $0\ne P\in F[[X]][Y]$ and put
$r=\ord_{(X,Y)}P$. There are indices $c_1,\ldots,c_r\in I$, with repetitions
permitted, such that
\[
 Q(X,Y)=\prod_{j=1}^r(Y-u_{c_j}(X))
 \quad\hbox{satisfies}\quad
 v(Q(X,u_b))\ge v(P(X,u_b))\quad\hbox{for every }b\in I.
\]
The empty product for $r=0$ is $1$.
\end{lemma}

\begin{proof}
Factor $P=a(X)\prod_{j=1}^D(Y-\alpha_j)$ in a valued algebraic extension of
$F((X))$, extending the $X$-adic valuation, and count roots with multiplicity.
The weighted Gauss valuation at weight one is multiplicative and gives
\[
 r=v(a)+\sum_{j=1}^D\min(1,v(\alpha_j)).
\]
This identity needs neither separability nor a Puiseux parametrization.
For completeness, the ordinary Gauss valuation is multiplicative because,
after dividing each polynomial by a coefficient of least valuation, its
nonzero reduction over the residue field has a nonzero product. Substituting
$Y=XZ$ proves the weighted version. Factoring into linear factors then gives
the displayed identity.

Let $A$ be the multiset of roots of valuation at least one, and write
$m=|A|$. The quantity
\[
 k=r-m=v(a)+\sum_{v(\alpha_j)<1}v(\alpha_j)
\]
is an integer. It is nonnegative: the ordinary Gauss valuation of $P$ is
nonnegative, and
\[
 k=G_0(P)+\sum_{0<v(\alpha_j)<1}v(\alpha_j)\ge0.
\]
For each $\alpha\in A$ choose a branch $u_{c(\alpha)}$ maximizing
$v(u_b-\alpha)$ over the finite set $B$. The ultrametric inequality implies,
for every $b$,
\[
 v(u_b-u_{c(\alpha)})\ge
 \min\{v(u_b-\alpha),v(u_{c(\alpha)}-\alpha)\}
 =v(u_b-\alpha).
\]
Every root of valuation below one contributes exactly its own valuation to
$v(P(X,u_b))$, since $v(u_b)\ge1$. Hence
\[
 Q_0=X^k\prod_{\alpha\in A}(Y-u_{c(\alpha)})
\]
dominates $P$ on all the branches and has total order $k+m=r$.
Replace each of its $k$ factors $X$ by $Y-u_{c_0}$, for any fixed $c_0\in I$.
This cannot decrease an evaluation valuation, because all contacts are at
least one. The resulting $Q$ has exactly $r$ branch factors.
\end{proof}

\begin{corollary}[Finite-record domination for formal series]\label{cor:formal}
If $H\notin\m^d$, where $\m=(X,Y)$, then a product of exactly $d-1$ branch
factors vanishes on every record on which $H$ vanishes.
\end{corollary}

\begin{proof}
Let $r=\ord_\m H<d$ and choose an integer $L>r$ with
$L\ge\max_b n_b$. The polynomial obtained by discarding terms of total degree
at least $L$ has total order $r$ and the same records as $H$, since
$u_b\in XF[[X]]$. Apply Lemma~\ref{lem:domination}. If $r<d-1$, append any
$d-1-r$ branch factors. Their evaluation valuations are nonnegative, so
previously vanishing records still vanish. The resulting product has total
order exactly $d-1$, as its coefficient of $Y^{d-1}$ is one.
\end{proof}

\begin{remark}
A stronger domination statement for the actual valuations of a formal series
follows by formal Weierstrass preparation after extracting its largest common
power of $X$. The finite-record statement above avoids needing that theorem
and is sufficient for all the distance and algorithmic claims here.
\end{remark}

\subsection{Exact distance, all errors and all erasures}

Define the target distance
\[
 D_d(\mathbf n;B)=\min_{H\notin\m^d}
       |\{b:\Omega_{\mathbf n}(H)_b\ne0\}|.
\]
For a vector $\mathbf m=(m_c)_{c\in I}$ of nonnegative integers, define
\[
 V_b(\mathbf m)=
 \begin{cases}
 +\infty,&m_b>0,\\
 \displaystyle\sum_{c\ne b}m_cs_{bc},&m_b=0.
 \end{cases}
\]
This convention avoids expressions such as $0\cdot\infty$.

\begin{theorem}[Exact unequal-record distance]\label{thm:distance}
For every $d\ge1$ and every nonnegative precision vector,
\[
 \boxed{\quad D_d(\mathbf n;B)=N-
 \max_{\substack{m_c\in\mathbb Z_{\ge0}\\\sum_cm_c=d-1}}
    |\{b:V_b(\mathbf m)\ge n_b\}|.\quad}
\]
Every maximum supplies a minimum-distance witness
$Q=\prod_c(Y-u_c)^{m_c}$. In particular the complete distance depends only on
the contact matrix and the precision vector, not on any other branch
coefficients. If $d>N$, the distance is zero at every finite precision.
\end{theorem}

\begin{proof}
Every displayed product has total order $d-1$ and its zero records are exactly
those satisfying the displayed inequalities. Conversely
Corollary~\ref{cor:formal} replaces any series with nonzero target by such a
product without reducing its zero set. Therefore the largest possible zero
set is exactly the finite maximum. If $d>N$, use each branch at least once
and pad with repeated factors to reach $d-1$ factors; this product vanishes
on all branches.
\end{proof}

\begin{corollary}[Sharp adversarial correction criterion]\label{cor:errors}
The target jet is recoverable for every $H$, despite at most $w$ arbitrarily
corrupted records and at most $z$ flagged erased records, if and only if
\[
 2w+z<D_d(\mathbf n;B).
\]
For a fixed erasure set $E$, delete its branches and use
$2w<D_d(\mathbf n|_{I\setminus E};B\setminus E)$ instead.
If all branches are erased, a positive-order target cannot be recovered.
\end{corollary}

\begin{proof}
Two explanations of a received word differ on at most the union of their
corruption sets and the erasure set, hence on at most $2w+z$ records. If their
target jets differed, their difference would contradict the distance bound.
Conversely take a minimum-distance product witness and the zero series.
Erase at most $z$ differing positions, and divide the remaining at most
$2w$ positions into two sets of size at most $w$. A word agreeing with one
explanation on one set and the other on the second has two admissible
explanations with different target jets. The distance-zero case requires no
errors at all.
\end{proof}

Nonnegative weights $a_b$ may be assigned to records. With weighted support
in place of cardinality the same proof gives
\[
 D_d^{(a)}=\sum_ba_b-
 \max_{\sum m_c=d-1}\sum_{V_b(\mathbf m)\ge n_b}a_b.
\]
For arbitrary indivisible weights the simple inequality $2w+z<D_d^{(a)}$ is
a sufficient correction condition; its necessity needs an actual partition
of differing weights into the permitted error budgets. We do not assert the
unweighted necessity proof for arbitrary weights.

\subsection{A contact-tree algorithm with explicit witnesses}

Represent the contacts by the usual rooted tree. For at least two leaves,
the root has incoming length $\ell_{\rm root}=\min_{b\ne c}s_{bc}$.
Children separate where the contact becomes larger; an internal node's
incoming length is its contact level minus its parent's level. The path
shared by two distinct leaves has total length $s_{bc}$. Terminal edges
may be assigned zero length; a mark at a leaf contributes infinity to that
leaf itself. Compress nodes with one child, leaving $O(N)$ nodes.

Let $B_v$ be the leaves below a node $v$. For $j\ge0$ and $t\ge0$ let
$A_v(j,t)$ be the largest total weight of leaves whose precision is met by
placing exactly $j$ marks in $B_v$, with repeats allowed, when contributions
from marks outside $B_v$ and from higher path segments already equal $t$ at
every leaf below $v$.

\begin{theorem}[Unequal-precision tree recursion]\label{thm:tree}
At a leaf $b$,
\[
 A_b(j,t)=\begin{cases}a_b,&j>0\text{ or }t\ge n_b,\\0,&\text{otherwise}.
 \end{cases}
\]
At an internal node with children $c$,
\[
 \boxed{\quad A_v(j,t)=
 \max_{\substack{j_c\ge0\\\sum_cj_c=j}}
       \sum_c A_c(j_c,t+j\ell_v).\quad}
\]
Consequently $D_d^{(a)}=\sum_ba_b-A_{\rm root}(d-1,0)$.
For integer contacts and precisions, the table through $j=d-1$ can be
computed in $O(Nd^2(n_{\max}+1))$ arithmetic operations, where
$n_{\max}=\max_b n_b$. Backtracking gives the mark multiplicities and thus
an explicit ambiguity witness.
\end{theorem}

\begin{proof}
All $j$ marks below an internal node share its incoming segment with every
leaf below that node. They therefore add $j\ell_v$ to the inherited value.
Once the marks are split among the children, no additional contribution
passes between distinct children; the remaining choices are independent.
At a leaf, a mark makes the polynomial vanish identically there, while an
unmarked leaf has precisely its inherited valuation. These facts prove the
recursion and its agreement with Theorem~\ref{thm:distance}.

Cap every state $t\ge n_{\max}$ at $n_{\max}$; all such leaves are already
covered. For each fixed child-input value $t'$, compute the max-convolution
\[
 C_v(j,t')=\max_{\sum j_c=j}\sum_cA_c(j_c,t')
\]
for all $j\le d-1$, then read
$A_v(j,t)=C_v(j,\min(n_{\max},t+j\ell_v))$.
The convolutions cost $O(\deg(v)d^2)$ at each $t'$, and the sum of node
degrees in a compressed tree is $O(N)$. This explains the stated bound
without an extra factor $d$ from recomputing convolutions for each state.
\end{proof}

The operation count is polynomial in the amount of observed data, and
pseudopolynomial if the lengths are supplied only in binary. The common
precision theory can also optimize by a greedy valuation sequence; unequal
leaf thresholds make the mark-allocation formulation useful.

\subsection{The unique minimum precision vector for maximum distance}

Suppose $1\le d\le N$. For each branch let $S_b(d-1)$ be the sum of its
$d-1$ largest contacts with other branches, with $S_b(0)=0$, and put
\[
 n_b^*=1+S_b(d-1).
\]

\begin{theorem}[Componentwise saturation and exact sampling cost]\label{thm:saturation}
For every precision vector,
\[
 D_d(\mathbf n;B)\le N-d+1,
\]
and equality holds if and only if $n_b\ge n_b^*$ for every $b$.
Thus $\mathbf n^*$ is the unique componentwise minimum precision vector
achieving maximum target distance. The smallest total number of recorded
field coefficients at maximum distance is exactly
\[
 \sum_bn_b^*=N+\sum_bS_b(d-1).
\]
The least common precision is $\max_bn_b^*$, so using unequal precision saves
exactly $N\max_bn_b^*-\sum_bn_b^*$ coefficients compared with the best
common-precision design.
\end{theorem}

\begin{proof}
A product of any $d-1$ distinct branch equations has total order $d-1$ and
vanishes on those $d-1$ branches, giving the upper bound.

If $n_b\le S_b(d-1)$, choose its $d-1$ highest-contact neighbours and multiply
their equations. The product vanishes on these neighbours and has valuation
$S_b(d-1)$ at $b$, so it has at least $d$ zero records and distance at most
$N-d$. This proves necessity, including $d=1$ when an empty record makes the
constant series $1$ invisible there.

For sufficiency assume all componentwise inequalities, and suppose that a
series $H$ has zero records on a $d$-element subset $S$. Divide $H$ by the
monic distinguished polynomial $\prod_{b\in S}(Y-u_b)$ to obtain a
remainder $R$ of $Y$-degree below $d$; the discarded multiple belongs to
$\m^d$. Lagrange interpolation over $F((X))$ gives
\[
 R(X,Y)=\sum_{b\in S}
 \frac{H(X,u_b)}{\prod_{c\in S\setminus\{b\}}(u_b-u_c)}
 \prod_{c\in S\setminus\{b\}}(Y-u_c).
\]
Each quotient has $X$-valuation at least
\[
 n_b-\sum_{c\in S\setminus\{b\}}s_{bc}\ge1,
\]
and each branch product has total order $d-1$. Hence $R\in\m^d$ and
$H\in\m^d$. No nonzero target can have $d$ zero records, so the distance
is at least $N-d+1$. The cost statements now follow componentwise.
\end{proof}

\begin{remark}
The sufficiency proof can avoid formal division by first truncating $H$ to
total degree at least $\max n_b$ and at least $d$, then using ordinary monic
polynomial division. Evaluation congruences and the target jet are
unchanged. Thus the formal machinery is not needed in the saturation proof
either.
\end{remark}

\paragraph{An asymmetric example.}
Over $F=\mathbb F_3$ take
\[
 B=(0,X^3,X,2X),\qquad d=2.
\]
Only the first two branches have contact $3$; all other contacts equal $1$.
The unique minimum saturation vector is
\[
 \mathbf n^*=(4,4,2,2).
\]
It records $12$ coefficients, compared with $16$ for the least common
precision, and has target distance $3$. Thus the three coefficients of
$J_2(H)$ are determined despite one arbitrarily corrupted branch record.
Reducing any one component of $\mathbf n^*$ destroys that guarantee.

More generally, over a field with at least $N-1$ elements take
$0,X^L$ and $N-2$ distinct nonzero scalar multiples of $X$, where $L>1$.
For $2\le d\le N$, the close pair has optimal precision $L+d-1$ and every
other branch has optimal precision $d$. The total is
$2L+Nd-2$, whereas common precision costs $N(L+d-1)$.
For fixed $N,d$ the ratio of common to unequal cost tends to $N/2$ as
$L\to\infty$.

\subsection{A weighted-jet extension}

Suppose, more specifically, that all $u_b\in X^\kappa F[[X]]$ for an integer
$\kappa\ge1$. Give $X$ weight one and $Y$ weight $\kappa$, and let
$J_{d,\kappa}(H)$ consist of the coefficients $[X^iY^j]H$ with
$i+\kappa j<d$. Put
\[
 k=\left\lfloor\frac{d-1}{\kappa}\right\rfloor,
 \qquad a=d-1-\kappa k,\qquad 0\le a<\kappa.
\]

\begin{corollary}[Exact distance for weighted jets]\label{cor:weighted}
Every ambiguity with nonzero $J_{d,\kappa}$ can be replaced, without losing
any zero record, by
\[
 X^a\prod_{c\in I}(Y-u_c)^{m_c},\qquad \sum_cm_c=k.
\]
Consequently its target distance is
\[
 D_{d,\kappa}(\mathbf n;B)=
 N-\max_{\sum m_c=k}|\{b:a+V_b(\mathbf m)\ge n_b\}|.
\]
The tree recursion starts with $k$ marks and inherited precision $a$.
If $k<N$, the maximum possible distance is $N-k$, and it is attained exactly
when
\[
 n_b\ge a+1+\text{the sum of the $k$ largest contacts $s_{bc}$, $c\ne b$}
 \quad\text{for every }b.
\]
If $k\ge N$, the distance is zero at every precision.
\end{corollary}

\begin{proof}
Use the weighted Gauss identity
$\ord_\kappa P=v(a_P)+\sum_\alpha\min(\kappa,v(\alpha))$
in the proof of Lemma~\ref{lem:domination}. Roots of valuation below
$\kappa$ contribute equally at every branch, while each remaining root may
be replaced by its closest branch. The resulting product of $m$ branch
factors and a power of $X$ has the same weighted order $r<d$ and dominates
the original polynomial. Multiply by $X^{d-1-r}$, and replace each group
of $\kappa$ residual factors $X$ by a branch factor. This cannot decrease
evaluation valuations, as $v(u_b-u_c)\ge\kappa$, and gives exactly the
displayed form with $k$ factors and residual exponent $a$. Truncating the
formal series at sufficiently high weighted degree justifies the same
argument for every finite collection of records.

For maximum distance, the product using $k$ distinct branches gives the
upper bound $N-k$. Necessity of each precision inequality follows by using
the $k$ closest neighbours of the branch in question and the factor $X^a$.
For sufficiency interpolate on any $k+1$ branches. The discarded full branch
product has weighted order $\kappa(k+1)\ge d$, and every Lagrange term has
weighted order at least $(a+1)+\kappa k=d$. Thus $k+1$ zero records force
the weighted target to vanish. If $k\ge N$, mark every branch and pad.
\end{proof}

\subsection{A decoder using two linear systems}

The exact distance criterion does not by itself give a polynomial-time
decoder at the entire optimal error radius. The following stronger,
computable condition does. It generalizes the error-locator linearization
behind Welch--Berlekamp decoding, while retaining the unrestricted higher
coefficients of the unknown series.

\begin{theorem}[A polynomial-time sufficient decoding condition]\label{thm:decode}
Assume arithmetic and linear algebra over $F$ are effective. Delete a known
erasure set, and retain the notation $B,\mathbf n$ for the surviving
branches and lengths. If at most $w$ remaining records are corrupted and
\[
 D_{d+w}(\mathbf n;B)>w,
\]
then $J_d(H)$ can be recovered by two finite linear systems, with no
enumeration of error sets and no factorization of an error locator. The
algorithm is polynomial in $N,n_{\max},d,w$ in field operations.
\end{theorem}

\begin{proof}
The condition implies $n_{\max}\ge d+w$: otherwise the nonzero-target series
$X^{n_{\max}}$ vanishes on every record at target order $d+w$.
Write the received record at branch $b$ as $y_b(X)$ modulo $X^{n_b}$.
Find $E,Q$ satisfying
\[
 \begin{split}
 E(X,Y)&=Y^w+\sum_{i=0}^{w-1}e_i(X)Y^i,\\
 &\hspace{-12pt}e_i(X)\in X^{w-i}F[X],\qquad
 Q\in F[X,Y],\quad\deg_{\rm total}Q<n_{\max},\\
 Q(X,u_b)&\equiv E(X,u_b)y_b(X)\pmod {X^{n_b}}\quad\text{for all }b.
 \end{split}
\]
The coefficient series $e_i$ may be truncated to terms with
$\deg_X(e_i)+i<n_{\max}$; their higher terms cannot affect any equation or
the subsequent target calculation. These are linear equations for the
unknown coefficients of $E$ and $Q$, since every $y_b$ is known, and the
coefficient of $Y^w$ in $E$ is fixed to one.

A solution exists. Let $C$ be the actual corruption set, of size at most
$w$, and take the product of $Y-u_b$ over $b\in C$, padded with arbitrary
branch factors to a monic product of degree $w$. This has the required
divisibilities. Let $Q$ be the total-degree truncation of $EH$ below
$n_{\max}$. The congruences hold on uncorrupted branches because $y_b$ is
correct, and on corrupted branches because $E$ vanishes there. Truncation
preserves the congruences. The actual locator is a formal polynomial in $Y$;
truncate its coefficients as specified above to obtain the finite $E$ in
the linear system. The omitted terms have total order at least
$n_{\max}$, so this truncation preserves both the congruences and the
required coefficient divisibilities.

Now use \emph{any} solution of the linear system, not necessarily the actual
error locator. On every uncorrupted branch the series $Q-EH$ has a zero
record, so it has at most $w$ nonzero records. By the assumed distance,
\[
 Q-EH\in\m^{d+w}.
\]
Since $E$ has total order exactly $w$, its initial homogeneous form is
nonzero in the domain $F[X,Y]$. Therefore multiplication by $E$ induces an
injective map
\[
 F[[X,Y]]/\m^d\longrightarrow
   \m^w/\m^{d+w}.
\]
Solve the second linear system
\[
 E h\equiv Q\pmod {\m^{d+w}},\qquad\deg_{\rm total}h<d.
\]
It is consistent, with $h=J_d(H)$, and injectivity makes its solution unique.

There are at most $\binom{n_{\max}+1}{2}+wn_{\max}$ unknowns in the first
system and $\sum_b n_b\le Nn_{\max}$ equations. The second has
$\binom{d+1}{2}$ unknowns and at most $\binom{d+w+1}{2}$ equations.
Ordinary Gaussian elimination therefore suffices with polynomial cost.
\end{proof}

\begin{remark}[Scope of the decoding condition]
The condition $D_{d+w}>w$ is generally stronger than the sharp uniqueness
condition $D_d>2w$. Indeed, appending $w$ branch factors to a
minimum-distance witness, at distinct currently nonzero records whenever
possible, proves $D_{d+w}\le\max(D_d-w,0)$. Thus the decoder condition
implies the uniqueness condition. At the asymmetric example's minimum vector
$(4,4,2,2)$, one-error uniqueness holds for $J_2$, but $D_3=0$ and the
linear-system theorem does not apply. Increasing the lengths to
$(5,5,3,3)$ gives $D_3=2>1$, so the theorem supplies a one-error decoder
using $16$ coefficients, compared with $20$ under its least common
precision. No assertion is made that this extra precision is necessary for
polynomial-time decoding.
\end{remark}

\subsection{Priority, limitations and further work}

The ingredients are classical. Bhargava's $P$-ordering theory gives regular
bases for integer-valued polynomials and ultrametric interpolation
\cite{Bhargava2009}. Amini gives a tree formulation of factorial sequences
\cite{Amini2016}. Root replacement by nearby integral points and a related
maximization recursion also appear in a public question on monic polynomial
root counts modulo prime powers \cite{MO}. These observations are antecedents
to the domination and tree method, rather than discoveries asserted here.
The error-locator construction is an adaptation of classical
interpolation-based decoding; see Gao \cite{Gao} for that circle of ideas.

Relative to the supplied article, the new statements are the unequal-record
distance formula for arbitrary formal-series jets, its weighted mark
algorithm, the exact componentwise saturation and total coefficient cost,
and the stated sufficient decoder. The common-length versions and the
exactly-$d$-branch clean interpolation threshold are already in that article.
The limited literature search did not establish global priority for the
combined statements. In particular this note should not be described as a
resolution of an established external named conjecture.

The results work in every characteristic, because the proofs do not divide
by factorials or require separability. They require a coefficient field.
Over $\mathcal O/\pi^r$ with $r>1$, nilpotent coefficient cancellation invalidates
the valued-field reduction, and the article already gives examples of extra
inverse-pole costs. Extending the exact distance to that setting, and finding
a polynomial-time decoder at the full sharp unequal-length radius, remain
open in this note. The mathematical proofs are human-readable proofs;
no claim of a compiled Lean formalization is made.

